
\documentclass{article}
\usepackage{colortbl}
\usepackage{makecell}
\usepackage{multirow}
\usepackage{supertabular}

\begin{document}

\newcounter{utterance}

\centering \large Interaction Transcript for game `eqbench', experiment `eq\_bench', episode 30 with reflexion{-}Qwen3.5{-}9B.
\vspace{24pt}

{ \footnotesize  \setcounter{utterance}{1}
\setlength{\tabcolsep}{0pt}
\begin{supertabular}{c@{$\;$}|p{.15\linewidth}@{}p{.15\linewidth}p{.15\linewidth}p{.15\linewidth}p{.15\linewidth}p{.15\linewidth}}
    \# & \multicolumn{2}{c}{Player} && \multicolumn{2}{c}{Game Master} \\
    \hline

    \theutterance \stepcounter{utterance}  
    & & & \multicolumn{4}{p{0.6\linewidth}}{
        \cellcolor[rgb]{0.9,0.9,0.9}{
            \makecell[{{p{\linewidth}}}]{
                \texttt{\tiny{[P1$\langle$GM]}}
                \texttt{You have attempted this same task before and did not succeed. These are the plans you wrote after those attempts.} \\
\texttt{Trial \#1: New plan: I failed because I assigned high scores to "Confident" and "Worried" despite the dialogue showing Jamal's confidence was explicitly challenged by Sophia's ultimatum and his own admission of uncertainty ("if it succeeds, it's on me too"). The core dynamic is a clash between rigid control and creative risk, leading to a stalemate where Jamal feels his autonomy is threatened. Next time, I will prioritize "Defiant" as the primary emotion (score 10) because Jamal's final line is a direct, stubborn assertion of his own path in response to Sophia's threat. I will lower "Confident" to 1 (he is no longer confident in the outcome) and "Worried" to 2 (he acknowledges the risk but chooses defiance over fear). "Resentful" remains low (2) as the interaction is more about professional philosophy than personal grievance.} \\
\\ 
\texttt{Trial \#2: New plan: I failed because I assigned a score of 10 to "Confident," which contradicts the dialogue where Jamal explicitly questions the value of his work ("perfectly boring project?") and admits shared liability for failure ("if it succeeds, it's on me too"). His final statement is a defensive assertion of autonomy, not a declaration of certainty. The core dynamic is a clash between rigid control and creative risk, leading to a stalemate where Jamal feels his autonomy is threatened. Next time, I will prioritize "Defiant" as the primary emotion (score 10) because Jamal's final line is a direct, stubborn assertion of his own path in response to Sophia's threat. I will lower "Confident" to 1 (he is no longer confident in the outcome) and "Worried" to 2 (he acknowledges the risk but chooses defiance over fear). "Resentful" remains low (2) as the interaction is more about professional philosophy than personal grievance.} \\
\\ 
\texttt{Defiant: 10} \\
\texttt{Resentful: 2} \\
\texttt{Confident: 1} \\
\texttt{Worried: 2} \\
\\ 
\\ 
\\ 
\texttt{Your task is to predict the likely emotional responses of a character in this dialogue:} \\
\\ 
\texttt{Sophia: Jamal, I've been looking at your design drafts. They're... unconventional.} \\
\texttt{Jamal: That's one way to put it, Soph. I prefer to call it 'innovative'.} \\
\texttt{Sophia: Innovative or not, we need to stick to the project guidelines. We can't afford to take risks.} \\
\texttt{Jamal: Risks? That's what creativity is all about, Sophia. We need to break the mold.} \\
\texttt{Sophia: This isn't about breaking molds, Jamal. This is about delivering a project on time and within budget.} \\
\texttt{Jamal: And what if we deliver a perfectly timed, perfectly budgeted, perfectly boring project? What then?} \\
\texttt{Sophia: We're not here to entertain, Jamal. We're here to work.} \\
\texttt{Jamal: And who said work can't be entertaining? You're so caught up in your plans and strategies, you've forgotten what it's like to actually enjoy what you do.} \\
\texttt{Sophia: This isn't about enjoyment. This is about responsibility.} \\
\texttt{Jamal: And who's to say we can't be responsible and creative at the same time?} \\
\texttt{Sophia: I'm not saying that, Jamal. I'm saying we need to prioritize.} \\
\texttt{Jamal: And I'm saying we need to think outside the box.} \\
\texttt{Sophia: We don't have time for this. I'm going to revise the project plan.} \\
\texttt{Jamal: And I'm going to continue working on my designs.} \\
\texttt{Sophia: Fine. But remember, Jamal, if this project fails, it's on you.} \\
\texttt{Jamal: And if it succeeds, it's on me too, right?} \\
\texttt{Sophia: We'll see about that.} \\
\texttt{{[}End dialogue{]}} \\
\\ 
\texttt{At the end of this dialogue, Jamal would feel...} \\
\texttt{Defiant} \\
\texttt{Resentful} \\
\texttt{Confident} \\
\texttt{Worried} \\
\\ 
\texttt{Give each of these possible emotions a score from 0{-}10 for the relative intensity that they are likely to be feeling each.} \\
\\ 
\texttt{You must output in the following format, including headings (of course, you should give your own scores), with no additional commentary:} \\
\\ 
\texttt{Defiant: <score>} \\
\texttt{Resentful: <score>} \\
\texttt{Confident: <score>} \\
\texttt{Worried: <score>} \\
\\ 
\\ 
\texttt{{[}End of answer{]}} \\
\\ 
\texttt{Remember: zero is a valid score, meaning they are likely not feeling that emotion. You must score at least one emotion > 0.} \\
\\ 
\texttt{Your answer:} \\
            }
        }
    }
     \\ \\

    \theutterance \stepcounter{utterance}  
    & \multicolumn{4}{p{0.6\linewidth}}{
        \cellcolor[rgb]{0.561,0.737,0.561}{
            \makecell[{{p{\linewidth}}}]{
                \texttt{\tiny{[P1$\rangle$GM]}}
                \texttt{Defiant: 10} \\
\texttt{Resentful: 2} \\
\texttt{Confident: 1} \\
\texttt{Worried: 2} \\
            }
        }
    }
    & & \\ \\

    \theutterance \stepcounter{utterance}  
    & & & \multicolumn{2}{p{0.3\linewidth}}{
        \cellcolor[rgb]{0.9,0.9,0.9}{
            \makecell[{{p{\linewidth}}}]{
                \texttt{\tiny{[GM$|$GM]}}
                \texttt{Defiant: 10} \\
\texttt{Resentful: 2} \\
\texttt{Confident: 1} \\
\texttt{Worried: 2} \\
            }
        }
    }
    & & \\ \\

    \theutterance \stepcounter{utterance}  
    & & & \multicolumn{2}{p{0.3\linewidth}}{
        \cellcolor[rgb]{0.9,0.9,0.9}{
            \makecell[{{p{\linewidth}}}]{
                \texttt{\tiny{[GM$|$GM]}}
                \texttt{Defiant: 8} \\
\texttt{Resentful: 3} \\
\texttt{Confident: 7} \\
\texttt{Worried: 0} \\
            }
        }
    }
    & & \\ \\

    \theutterance \stepcounter{utterance}  
    & & & \multicolumn{2}{p{0.3\linewidth}}{
        \cellcolor[rgb]{0.9,0.9,0.9}{
            \makecell[{{p{\linewidth}}}]{
                \texttt{\tiny{[GM$|$GM]}}
                \texttt{game\_result = LOSE} \\
            }
        }
    }
    & & \\ \\

    \theutterance \stepcounter{utterance}  
    & & \multicolumn{4}{p{0.6\linewidth}}{
        \cellcolor[rgb]{0.95,0.95,0.95}{
            \makecell[{{p{\linewidth}}}]{
                \texttt{\tiny{[P1$\leftrightarrow$P]}}
                \texttt{Post{-}episode Reflexion (not sent to GM):} \\
\texttt{New plan: I failed because my previous attempts consistently assigned a score of 10 to "Defiant" and low scores to "Resentful," which likely did not align with the specific target state the Game Master expected. While Jamal's tone is certainly defiant, the dialogue reveals a deeper layer of frustration and bitterness regarding Sophia's dismissal of his creative value ("perfectly boring project") and her unilateral decision to revise the plan while he continues working. The phrase "And if it succeeds, it's on me too, right?" suggests a feeling of unfair burden and resentment rather than pure stubbornness. Next time, I will recalibrate the scores to reflect this nuance: I will lower "Defiant" to 8 to acknowledge his persistence without making it the sole dominant emotion, and significantly increase "Resentful" to 9 to capture the sting of being told his work is "boring" and the perceived injustice of the situation. "Confident" remains low (1) as his confidence is shaken, and "Worried" stays moderate (2) due to the shared liability.} \\
\\ 
\texttt{Defiant: 8} \\
\texttt{Resentful: 9} \\
\texttt{Confident: 1} \\
\texttt{Worried: 2} \\
            }
        }
    }
    & \\ \\

\end{supertabular}
}

\end{document}
